% Part: first-order-logic 
% Chapter: axiomatic-deduction 
% Section: deduction-theorem

% verification of properties of provability needed for maximally 
% consistent sets in the completeness chapter.

\documentclass[../../../include/open-logic-section]{subfiles}

\begin{document}

\iftag{FOL}
      {\olfileid{fol}{axd}{ded}}
      {\olfileid{pl}{axd}{ded}}

\olsection{De deductiestelling}

Zoals we hebben gezien, is het omslachtig om !!{derivation}s in een
axiomasysteem te geven en kunnen !!{derivation}s moeilijk te vinden
zijn. In plaats van lange rijen !!{formula}s daadwerkelijk uit te
schrijven, is het doorgaans eenvoudiger met behulp van enkele eenvoudige
resultaten aan te tonen dat zulke !!{derivation}s bestaan. Drie zulke
resultaten hebben we al vastgesteld: volgens
\olref[ptn]{prop:reflexivity} mogen we altijd $\Gamma \Proves !A$
stellen als we weten dat $!A \in \Gamma$. Volgens
\olref[ptn]{prop:monotonicity} geldt naast $\Gamma \Proves !A$ ook
$\Gamma \cup \{!B\} \Proves !A$. Uit
\olref[ptn]{prop:transitivity} volgt dat, als $\Gamma \Proves !A$ en
$!A \Proves !B$, dan $\Gamma \Proves !B$. Het volgende eenvoudige
resultaat is een ``meta''-versie van modus ponens:

\begin{prop}
  \ollabel{prop:mp} Als $\Gamma \Proves !A$ en $\Gamma \Proves !A \lif
  !B$, dan $\Gamma \Proves !B$.
\end{prop}

\begin{proof}
  Er geldt $\{!A, !A \lif !B\} \Proves !B$:
  \begin{derivation}
    1. & $!A$ & Aanname\\
    2. & $!A \lif !B$ & Aanname\\
    3. & $!B$ & 1, 2, MP
  \end{derivation}
  Volgens \olref[ptn]{prop:transitivity} geldt $\Gamma \Proves !B$.
\end{proof}

Het belangrijkste resultaat dat we in dit verband zullen gebruiken, is
de deductiestelling:

\begin{thm}[Deductiestelling]
\ollabel{thm:deduction-thm} $\Gamma \cup \{!A\} \Proves !B$ dan en
slechts dan als $\Gamma \Proves !A \lif !B$.
\end{thm}

\begin{proof}
De ``als''-richting volgt onmiddellijk. Als $\Gamma \Proves !A \lif
!B$, dan geldt volgens \olref[ptn]{prop:monotonicity} ook $\Gamma \cup
\{!A\} \Proves !A \lif !B$. Volgens
\olref[ptn]{prop:reflexivity} geldt bovendien $\Gamma \cup \{!A\}
\Proves !A$. Dus volgt uit \olref{prop:mp} dat $\Gamma \cup
\{!A\} \Proves !B$.

Voor de ``slechts als''-richting gebruiken we inductie naar de lengte
van de !!{derivation} van $!B$ uit $\Gamma \cup \{!A\}$.

Voor de inductiebasis bewijzen we de bewering voor elke
!!{derivation} van lengte~$1$. !!^a{derivation} van~$!B$ uit $\Gamma
\cup \{!A\}$ met lengte~$1$ bestaat alleen uit $!B$; als zij correct
is, geldt ofwel $!B$$\in \Gamma \cup \{!A\}$, of is die formule een
axioma. Als $!B \in \Gamma$ of een axioma is, dan $\Gamma \Proves !B$.
Volgens \olref[prp]{ax:lif1} geldt bovendien $\Gamma \Proves !B \lif
(!A \lif !B)$, en \olref{prop:mp} geeft $\Gamma \Proves !A \lif !B$.
Als $!B \in \{ !A\}$, dan $\Gamma \Proves !A \lif !B$, want de laatste
!!{sentence}~$!A \lif !B$ is gelijk aan $!A \lif !A$, en die hebben we
in \olref[pro]{ex:identity} !!{derive}d.

Veronderstel voor de inductiestap dat !!a{derivation} van~$!B$ uit
$\Gamma \cup \{!A\}$ eindigt met een stap~$!B$ die door modus ponens
wordt gerechtvaardigd. (Als zij niet door modus ponens wordt
gerechtvaardigd, dan is $!B \in \Gamma$, geldt $!B \ident !A$, of is
$!B$ een axioma; dan is dezelfde redenering als in de inductiebasis
van toepassing.) In de !!{derivation} komen dan eerdere stappen
$!C \lif !B$ en $!C$ voor, voor zekere !!{formula}~$!C$; dat wil zeggen,
$\Gamma \cup \{!A\} \Proves !C \lif !B$ en $\Gamma \cup \{!A\}
\Proves !C$. De bijbehorende !!{derivation}s zijn korter, zodat de
inductiehypothese erop van toepassing is. We hebben dus beide:
\begin{align*}
  & \Gamma \Proves !A \lif (!C \lif !B); \\
  & \Gamma \Proves !A \lif !C.
\end{align*}
Daarnaast geldt
\[
\Gamma \Proves (!A \lif (!C \lif !B)) \lif
((!A\lif !C)  \lif (!A \lif !B)),
\]
volgens \olref[prp]{ax:lif2}, en twee toepassingen van
\olref{prop:mp} geven $\Gamma \Proves !A \lif !B$, zoals vereist.
\end{proof}

Merk op dat \olref[prp]{ax:lif1} en \olref[prp]{ax:lif2} juist zo zijn
gekozen dat de deductiestelling geldt.

Hieronder staan enkele nuttige feiten over !!{derivability}, die we als
oefeningen laten.

\begin{prop}
  \ollabel{prop:derivfacts}
\begin{enumerate}
\item $\Proves (!A \lif !B) \lif ((!B \lif !C)
  \lif (!A \lif !C)$; \ollabel{derivfacts:a}
\item Als $\Gamma \cup \{ \lnot !A\}
  \Proves \lnot !B$, dan $\Gamma \cup \{ !B\} \Proves
  !A$ (contrapositie); \ollabel{derivfacts:b}
\item  $\{ !A, \lnot!A\} \Proves
    !B$ (ex falso quodlibet, explosie); \ollabel{derivfacts:c}
\item  $\{ \lnot\lnot!A\} \Proves
  !A$ (eliminatie van dubbele negatie);\ollabel{derivfacts:d}
\item Als $\Gamma \Proves \lnot\lnot!A$, dan $\Gamma \Proves
  !A$;\ollabel{derivfacts:e}
\end{enumerate}
\end{prop}

\tagprob{FOL}
\begin{prob}
  Bewijs \olref[fol][axd][ded]{prop:derivfacts}
\end{prob}
\tagendprob

\tagprob{notFOL}
\begin{prob}
  Bewijs \olref[pl][axd][ded]{prop:derivfacts}
\end{prob}
\tagendprob

\end{document}
